% ---
% id: foundations/cor:least-eigenvalue-over-enclosure
% kind: cor
% title: "Least eigenvalue and rank over an enclosure"
% summary: >-
%   A Hermitian matrix within ρ of a Hermitian M in Frobenius norm has its
%   spectrum within ρ of M's extreme eigenvalues. An entrywise enclosure whose
%   exactly zero rows are dropped and whose rest certifies a positive least
%   eigenvalue fixes the rank and support of every Hermitian matrix it contains.
% prominence: technical
% status: implemented
% uses: [lem:least-eigenvalue-bound, def:exact-reference]
% mentions: []
% assumes: []
% refs: []
% provenance: newly-derived
% implements: certiqs_key.foundations.dyadic:enclosure_eigenvalue_bounds
% lean:
% ---

Lemma~\ref{lem:least-eigenvalue-bound} bounds the least eigenvalue of a
matrix known exactly. An evaluation in ball arithmetic knows a matrix only
through an enclosure: part~(i) carries the bound across a ball in Frobenius
norm around an exact matrix, which is Weyl's inequality at the two ends of the
spectrum, and part~(ii) reads an \tref{entrywise enclosure} as such a ball.
Parts~(iii) and~(iv) decide the rank of every matrix an enclosure contains at
once: the rows the enclosure fixes at zero are removed, and the rest must be
positive definite throughout the enclosure. Where the rest is not shown to be,
and in particular where its centre vanishes, the enclosure decides no rank,
and the corollary claims none.

\begin{corollary}\label{cor:least-eigenvalue-over-enclosure}
Let $n\ge1$, and for a complex $n\times n$ matrix $X$ write
$\|X\|_{\mathrm F}\coloneqq\big(\sum_{i,l}|X_{il}|^{2}\big)^{1/2}$ and
$\opn{X}\coloneqq\max\{\|Xx\|:x\in\mathbb{C}^{n},\ \|x\|=1\}$.
\begin{enumerate}[label=(\roman*),leftmargin=2.4em]
\item \emph{(A ball in Frobenius norm.)} Let $R$ and $M$ be Hermitian
  $n\times n$ matrices and $\rho\ge0$ with $\|R-M\|_{\mathrm F}\le\rho$. Then
  $\opn{R-M}\le\|R-M\|_{\mathrm F}$, and
  \begin{equation}\label{eq:frobenius-ball}
    \lmin(R)\;\ge\;\lmin(M)-\rho,
    \qquad
    \lmax(R)\;\le\;\lmax(M)+\rho,
  \end{equation}
  so that every eigenvalue of $R$ lies in
  $[\lmin(M)-\rho,\lmax(M)+\rho]$. In particular $\lmin(R)\ge\ell-\rho$ for
  every lower bound $\ell$ on $\lmin(M)$, the right-hand side of
  \eqref{eq:least-eigenvalue-bound} for $M$ among them.
\item \emph{(An entrywise enclosure.)} If $(C,r)$ is an entrywise enclosure of
  $R$ (Definition~\ref{def:exact-reference}) and $\rho$ is a rational with
  $\rho\ge0$ and $\rho^{2}\ge\sum_{i,l}r_{il}^{2}$, then
  $\|R-C\|_{\mathrm F}\le\rho$ and (i) applies with $M=C$; whether $\rho$
  qualifies is a comparison of rationals.
\item \emph{(A factor form fixes the rank.)} Let $1\le k\le n$, let $E$ be a
  $k\times n$ matrix of rank $k$ and $\widetilde X$ a positive definite
  $k\times k$ matrix, and put $X\coloneqq E^{\dagger}\widetilde XE$. Then
  $X\succeq0$, $\operatorname{ran}X=\operatorname{ran}E^{\dagger}$ and
  $\operatorname{rank}X=k$.
\item \emph{(The rank over an entrywise enclosure.)} Let $C$ be a Hermitian
  $n\times n$ matrix with entries in $\mathbb{Q}(i)$ and $r_{il}=r_{li}\ge0$
  rationals, and say that $(C,r)$ \emph{contains} a Hermitian $X$ when
  $|X_{il}-C_{il}|\le r_{il}$ for all $i,l$. Let
  $Z\coloneqq\{i:C_{il}=0\text{ and }r_{il}=0\text{ for every }l\}$, the rows
  the pair fixes at zero, let $K$ be its complement in $\{1,\dots,n\}$ and
  $k\coloneqq|K|$. If $k=0$, the zero matrix is the only Hermitian matrix
  $(C,r)$ contains. If $k\ge1$, let $E$ be the $k\times n$ matrix whose rows
  are the transposed standard basis vectors $e_i^{\transpose}$, $i\in K$, in
  increasing order; then every Hermitian $X$ that $(C,r)$ contains satisfies
  $X=E^{\dagger}(EXE^{\dagger})E$. If moreover $\ell_K\le\lmin(ECE^{\dagger})$
  and a rational $\rho_K\ge0$ with $\rho_K^{2}\ge\sum_{i,l\in K}r_{il}^{2}$
  satisfy $\ell_K-\rho_K>0$, then for every Hermitian $X$ that $(C,r)$
  contains, $EXE^{\dagger}\succ0$, and so, by (iii), $X\succeq0$,
  $\operatorname{ran}X=\operatorname{span}\{e_i:i\in K\}$ and
  $\operatorname{rank}X=k$: the factor $E$, the support and the rank are the
  same for every such $X$.
\end{enumerate}
\end{corollary}

\begin{proof}
\emph{(i).} Let $A$ be a complex $n\times n$ matrix and $x$ a unit vector.
By the Cauchy--Schwarz inequality applied to each row,
$\|Ax\|^{2}=\sum_i\big|\sum_lA_{il}x_l\big|^{2}
\le\sum_i\big(\sum_l|A_{il}|^{2}\big)\|x\|^{2}=\|A\|_{\mathrm F}^{2}$, so
$\opn{A}\le\|A\|_{\mathrm F}$, and $|x^{\dagger}Ax|\le\|x\|\,\|Ax\|
\le\|A\|_{\mathrm F}$. With $A=R-M$, for every unit $x$,
$x^{\dagger}Rx=x^{\dagger}Mx+x^{\dagger}(R-M)x$ lies between
$x^{\dagger}Mx-\rho$ and $x^{\dagger}Mx+\rho$, hence between
$\lmin(M)-\rho$ and $\lmax(M)+\rho$, since the Rayleigh quotient of the
Hermitian $M$ lies between its extreme eigenvalues. The least and the
largest eigenvalue of the Hermitian $R$ are the minimum and the maximum of
$x^{\dagger}Rx$ over unit $x$, which gives \eqref{eq:frobenius-ball}, and
every eigenvalue of $R$ lies between them. The last sentence follows from
$\lmin(M)\ge\ell$.

\emph{(ii).} By \eqref{eq:entrywise-enclosure},
$\|R-C\|_{\mathrm F}^{2}=\sum_{i,l}|R_{il}-C_{il}|^{2}\le\sum_{i,l}r_{il}^{2}
\le\rho^{2}$, and $C$ is Hermitian by Definition~\ref{def:exact-reference}.
The sum of the squares of finitely many rationals is rational.

\emph{(iii).} For every $x$,
$x^{\dagger}Xx=(Ex)^{\dagger}\widetilde X(Ex)\ge0$, so $X\succeq0$, and
$x^{\dagger}Xx=0$ exactly when $Ex=0$, $\widetilde X$ being positive
definite. For a positive semidefinite $X$, $Xx=0$ exactly when
$x^{\dagger}Xx=0$, since $x^{\dagger}Xx=\|X^{1/2}x\|^{2}$ and
$Xx=X^{1/2}(X^{1/2}x)$. So $\ker X=\ker E$, and, $X$ being Hermitian,
$\operatorname{ran}X=(\ker X)^{\perp}=(\ker E)^{\perp}
=\operatorname{ran}E^{\dagger}$, whose dimension is
$\operatorname{rank}E=k$.

\emph{(iv).} Let $X$ be Hermitian and contained in $(C,r)$. For $i\in Z$ and
every $l$, $|X_{il}|=|X_{il}-C_{il}|\le r_{il}=0$, so row $i$ of $X$
vanishes, and so does column $i$, $X$ being Hermitian. If $k=0$, every row
vanishes and $X=0$, which $(C,r)$ contains. If $k\ge1$, $E^{\dagger}E$ is the
diagonal matrix with entry $1$ at the indices of $K$ and $0$ at those of $Z$,
so $E^{\dagger}EXE^{\dagger}E$ is $X$ with the rows and columns of $Z$ set
to zero, which is $X$. The matrix $EXE^{\dagger}$ is the principal
submatrix of $X$ on $K$, Hermitian, and $|(EXE^{\dagger})_{il}
-(ECE^{\dagger})_{il}|\le r_{il}$ for $i,l\in K$, so
$\|EXE^{\dagger}-ECE^{\dagger}\|_{\mathrm F}\le\rho_K$ as in (ii), with
$ECE^{\dagger}$ Hermitian. By (i),
$\lmin(EXE^{\dagger})\ge\ell_K-\rho_K>0$, so $EXE^{\dagger}$ is positive
definite, and (iii) applies with $\widetilde X=EXE^{\dagger}$; $E$ has rank
$k$, its rows being distinct standard basis vectors, and
$\operatorname{ran}E^{\dagger}=\operatorname{span}\{e_i:i\in K\}$. Neither
$E$ nor $k$ depends on $X$.
\end{proof}
